% Figure: hydroelectric penstock layout. Reservoir, surge tank, sloping
% penstock, turbine valve at the foot. The geometry the water-hammer case
% study carries numbers through.
\begin{figure}[htbp]
\centering
\begin{tikzpicture}[
  wall/.style={draw=engblack, very thick},
  pipe/.style={draw=engblue, very thick},
  water/.style={fill=engblue!15},
  flow/.style={draw=engblue, -{Stealth}, thick},
  dim/.style={draw=enggray, {Stealth}-{Stealth}, thin},
  label/.style={font=\scriptsize},
]
% Reservoir block (top left)
\fill[water] (-5,3.2) rectangle (-2.2,4.4);
\draw[wall] (-5,4.4) -- (-5,3.2) -- (-2.2,3.2);
\draw[engblue, thick] (-5,4.4) -- (-2.2,4.4);
\node[label, engblue] at (-3.6,4.0) {reservoir};
\node[label] at (-3.6,4.62) {free surface};
% Headrace from reservoir to surge tank
\draw[pipe] (-2.2,3.5) -- (-0.6,3.5);
% Surge tank (vertical riser)
\draw[wall] (-0.7,3.2) -- (-0.7,5.2);
\draw[wall] (-0.5,3.2) -- (-0.5,5.2);
\fill[water] (-0.7,3.2) rectangle (-0.5,4.7);
\node[label, anchor=west] at (-0.45,5.0) {surge tank};
% Penstock: sloping pipe down to the valve
\draw[pipe] (-0.6,3.35) -- (4.0,-1.2);
\draw[pipe] (-0.6,3.05) -- (4.2,-1.45);
% Flow arrow along penstock
\draw[flow] (0.6,2.4) -- (2.2,0.85);
\node[label, engblue] at (1.8,1.9) {$u_0$};
% Valve at the foot
\draw[wall, fill=engamber!30] (3.95,-1.25) rectangle (4.35,-1.0);
\node[label, anchor=west] at (4.4,-1.1) {valve};
% Turbine downstream of valve
\draw[wall] (4.6,-1.6) -- (5.3,-1.6) -- (5.3,-0.7) -- (4.6,-0.7) -- cycle;
\node[label] at (4.95,-1.15) {turbine};
% Datum line at the foot
\draw[enggray, dashed] (-5.2,-1.5) -- (5.6,-1.5);
\node[label, enggray, anchor=west] at (5.0,-1.7) {datum};
% Gross-head dimension
\draw[dim] (-5.6,-1.5) -- (-5.6,4.4);
\node[label, anchor=east] at (-5.65,1.5) {gross head $H$};
% Penstock length label
\node[label, engblue, rotate=-44] at (2.0,1.2) {length $L$};
\end{tikzpicture}
\caption[Hydroelectric penstock layout for the water-hammer case
study]{Layout of a hydroelectric scheme: a reservoir feeds a sloping
penstock of length $L$ that drops a gross head $H$ to a valve and
turbine at the foot. A surge tank near the top of the penstock gives
the water column a free surface to vent into when the valve closes,
limiting the pressure rise. The water-hammer case study carries
numbers through this geometry: steady flow at speed $u_0$, sudden
valve closure, the Joukowsky pressure rise, and the surge-tank or
slow-closure mitigation that keeps the peak pressure inside the pipe's
rating.}
\label{fig:vol03ch11:penstock-layout}
\end{figure}
